\chapter{Comunidades com Mais Relações}\label{ap:comunidades}

Este apêndice reproduz os resumos das dez comunidades de nível 1 com maior número de relações no índice construído (Seção~\ref{sec:met-hierarquia}). Os resumos foram gerados pelo LLM na etapa de sumarização de comunidades do GraphRAG e constituem o material consultado pelos modos \textit{global} e \textit{drift}; o título e o resumo de cada comunidade são apresentados tal como gerados, exceto pela substituição de nomes de pessoas pelo marcador [nome omitido], feita para preservar a confidencialidade dos dados. A Tabela~\ref{tab:ap-comunidades} resume o tamanho dessas comunidades.

\begin{table}[H]
\centering
\caption[Comunidades de nível 1 com mais relações]{Dez comunidades de nível 1 com maior número de relações no índice.}
\label{tab:ap-comunidades}
\small
\begin{tabular}{crr}
\toprule
Comunidade & Entidades & Relações \\
\midrule
169 & 3.980 & 4.216 \\
331 & 902 & 1.826 \\
329 & 683 & 1.203 \\
377 & 397 & 931 \\
382 & 476 & 915 \\
378 & 466 & 881 \\
330 & 418 & 818 \\
328 & 420 & 761 \\
46 & 523 & 749 \\
459 & 419 & 725 \\
\bottomrule
\end{tabular}
\par\smallskip{\footnotesize Fonte: elaborado pelo autor a partir do índice construído.}
\end{table}

\subsection*{Comunidade 169: Rede de Cuidado Hepatológico e Comorbidades Sistêmicas: PACIENTE, HEPATOLOGISTA e Diagnósticos Laboratoriais}

{\small 3.980 entidades e 4.216 relações.}

A comunidade analisada estrutura-se em torno do constructo do PACIENTE, que atua como nodo central em um ecossistema assistencial complexo e multifacetado. As entidades interconectadas abrangem desde marcadores laboratoriais críticos (FUNÇÃO RENAL, ANTI-HIV) e procedimentos históricos (TRANSFUSÃO SANGUÍNEA, HISTERECTOMIA, HERNIORRAFIA) até condições crônicas (INSUFICIÊNCIA RENAL CRÔNICA, CÂNCER DE PELE) e a figura coordenadora do HEPATOLOGISTA. As relações mapeadas revelam um fluxo contínuo de monitoramento diagnóstico, estratificação de risco e necessidade de integração interdisciplinar, evidenciando como a hepatologia contemporânea opera na interseção entre vigilância virológica, manejo de comorbidades e determinantes sociais da saúde.

\subsection*{Comunidade 331: Rede de Cuidado Hepatológico: Paciente, MASLD e Monitoramento de Biomarcadores}

{\small 902 entidades e 1.826 relações.}

A comunidade analisada estrutura-se em torno de um paciente central conectado a uma rede complexa de entidades diagnósticas, comorbidades metabólicas, fatores de risco comportamentais e intervenções farmacológicas. As relações mapeadas revelam um fluxo assistencial longitudinal, no qual biomarcadores laboratoriais (TGO, TGP, GGT, fosfatase alcalina, coagulação) e exames de imagem (elastografia, TC de abdome) são utilizados para monitorar a progressão de hepatopatias crônicas, incluindo MASLD/METALD, hepatopatia por esquistossomose e doença esteatótica não alcoólica. A presença de comorbidades como DM2 descompensada, pré-diabetes e sobrepeso, associada a padrões alimentares desequilibrados e histórico de uso de medicamentos hepatotóxicos, evidencia a necessidade de modelos de cuidado colaborativo. A integração desses dados permite estratificar riscos, ajustar condutas terapêuticas e otimizar a comunicação entre especialistas, laboratórios e o paciente.

\subsection*{Comunidade 329: Rede de Cuidado Hepatológico: Paciente, Santa Casa e Biomarcadores Tumorais}

{\small 683 entidades e 1.203 relações.}

A comunidade analisada estrutura-se em torno de um eixo central de cuidado hepatológico longitudinal, onde o paciente atua como nó integrador de uma rede multidisciplinar complexa. As entidades relacionam-se por meio de fluxos diagnósticos sequenciais (ultrassonografia, tomografia e ressonância magnética), monitoramento laboratorial contínuo (AFP, CEA e painel de marcadores tumorais) e intervenções cirúrgicas e nutricionais coordenadas por uma instituição de referência (Santa Casa). A dinâmica revela um modelo de atenção focado na vigilância de hepatopatias crônicas e neoplasias, com forte ênfase na correlação multimodal de exames, estratificação de risco baseada em histórico familiar e gestão integrada de comorbidades metabólicas e sistêmicas.

\subsection*{Comunidade 377: Rede de Cuidado em Diabetes Mellitus Tipo 2 e Hepatopatias Crônicas}

{\small 397 entidades e 931 relações.}

A comunidade analisada estrutura-se em torno do manejo longitudinal do Diabetes Mellitus Tipo 2 (DM2) e suas intersecções com hepatopatias crônicas e complicações sistêmicas. Os atores principais incluem pacientes, endocrinologistas, clínicos gerais e laboratórios de diagnóstico, conectados por fluxos de exames seriados (HbA1c, glicemia de jejum, microalbuminúria) e esquemas terapêuticos complexos (insulinas, metformina, pioglitazona). A dinâmica revela uma evolução do cuidado que transita entre fases compensadas e descompensadas, influenciada por determinantes sociais, histórico familiar e adesão medicamentosa. A integração entre dados laboratoriais e encaminhamentos especializados demonstra um modelo de cuidado colaborativo, ainda que com lacunas na comunicação e no rastreamento precoce de comorbidades hepáticas e renais.

\subsection*{Comunidade 382: Rede de Cuidado Hepatometabólico e Nutricional do Paciente}

{\small 476 entidades e 915 relações.}

A comunidade analisada estrutura-se em torno de um eixo central de cuidado multidisciplinar voltado ao manejo da doença hepática metabólica e da síndrome metabólica. As entidades principais incluem profissionais de nutrição, intervenções de atividade física, monitoramento laboratorial (como ALT e sorologias) e comorbidades associadas. Os fluxos relacionais demonstram uma rede de referência e contrarreferência focada na modificação do estilo de vida, controle ponderal e acompanhamento especializado, embora evidencie barreiras de adesão, lacunas no acesso a medicamentos e procedimentos cirúrgicos, e a necessidade de integração contínua entre dados clínicos e determinantes sociais.

\subsection*{Comunidade 378: Gestão Integrada de HAS e Hepatopatia Crônica: [nome omitido] e [nome omitido]}

{\small 466 entidades e 881 relações.}

A comunidade analisada estrutura-se em torno de um paciente central com múltiplas comorbidades crônicas, destacando-se a Hipertensão Arterial Sistêmica (HAS) e condições hepáticas avançadas. A rede conecta entidades farmacológicas (anti-hipertensivos, diuréticos), achados clínicos (PA elevada, AVC isquêmico, cirrose com ascite), fluxos diagnósticos (exames laboratoriais e de função renal) e atores assistenciais (médico [nome omitido] e cuidadora [nome omitido]). As relações evidenciam um modelo de cuidado longitudinal marcado por polifarmácia, necessidade de reconciliação medicamentosa, monitoramento renal e integração entre atenção cardiovascular e hepatológica, com lacunas na adesão terapêutica e na percepção do paciente sobre seu histórico familiar e controle pressórico.

\subsection*{Comunidade 330: Paciente, Fígado e Rede de Diagnóstico por Imagem Abdominal}

{\small 418 entidades e 818 relações.}

A comunidade analisada estrutura-se em torno de um eixo central de monitoramento hepatorrenal longitudinal, onde o paciente e o fígado funcionam como nós principais conectados a uma rede diversificada de exames de imagem (ultrassonografia, tomografia e ressonância magnética) e achados clínicos. As entidades revelam uma trajetória de cuidado focada na evolução de hepatopatias crônicas, comorbidades renais (cálculos, cistos, doença policística) e achados incidentais, demonstrando como a integração diagnóstica sustenta fluxos assistenciais multidisciplinares. A dinâmica observada destaca a importância da correlação imagiológica sequencial para estratificação de risco, ajuste terapêutico e coordenação entre especialidades, refletindo um modelo de cuidado centrado na evolução clínica e na otimização de desfechos.

\subsection*{Comunidade 328: Rede de Cuidado Hepatológico: HU, [nome omitido] e Fluxos Diagnósticos por Imagem}

{\small 420 entidades e 761 relações.}

A comunidade analisada estrutura-se em torno de um eixo central de diagnóstico por imagem e acompanhamento clínico longitudinal de hepatopatias crônicas. O Hospital Universitário atua como polo coordenador do cuidado, integrando pacientes como [nome omitido] a uma rede multidisciplinar. Exames como Colangio-RM, CPRM e Ressonância Magnética de Abdome funcionam como nós críticos para mapear a progressão de doenças (esteatose, cirrose, sobrecarga de ferro, síndrome colestática e Doença de Caroli), orientando condutas terapêuticas e fluxos de referência. A análise revela uma dinâmica assistencial complexa, onde a integração de dados laboratoriais, de imagem e clínicos é essencial para a estratificação de risco e a otimização de desfechos.

\subsection*{Comunidade 46: Paciente e Rede de Monitoramento Hepático-Metabólico}

{\small 523 entidades e 749 relações.}

A comunidade analisada estrutura-se em torno de um paciente central, interligado a uma extensa rede de parâmetros laboratoriais, diagnósticos clínicos e fluxos de acompanhamento longitudinal. As entidades representam marcadores hematológicos, metabólicos, imunológicos e hepatológicos que, em conjunto, delineiam um perfil clínico complexo, caracterizado por hepatopatia crônica de etiologia mista, comorbidades endócrinas e reumatológicas, e vigilância contínua de desequilíbrios eletrolíticos e proteicos. As relações mapeadas evidenciam um modelo de cuidado baseado em dados, onde a interpretação seriada de exames direciona ajustes terapêuticos, exclusão de hipóteses diagnósticas e integração entre especialidades, refletindo a dinâmica de um sistema de saúde focado na evolução clínica e na precisão diagnóstica.

\subsection*{Comunidade 459: Hepatopatia Crônica, Cirrose e Manejo de Ascite/Edema: Rede de Cuidado Multidisciplinar}

{\small 419 entidades e 725 relações.}

A comunidade analisada estrutura-se em torno de pacientes com hepatopatia crônica e cirrose hepática (alcoólica e viral B), cujas trajetórias clínicas são marcadas por complicações descompensadas como ascite grau 3, edema de membros inferiores e distúrbios eletrolíticos. A rede conecta sintomas físicos, exames laboratoriais (TAP, sódio, potássio), métodos de imagem (RNM abdominal) e intervenções farmacológicas (furosemida/Lasix, espironolactona), demonstrando um fluxo assistencial que depende da integração entre diagnóstico, monitoramento contínuo e ajuste terapêutico. A presença de comorbidades cardiovasculares e renais, aliada a fatores socioambientais, reforça a necessidade de modelos colaborativos para otimizar desfechos e prevenir descompensações graves.
